% ══════════════════════════════════════════════════════════════════════════════
%  CHAPITRE 4 — RÉALISATION
% ══════════════════════════════════════════════════════════════════════════════

\section{Introduction}

Ce chapitre décrit la réalisation concrète de la plateforme : l'environnement de développement utilisé, les choix technologiques, les interfaces développées, ainsi que les tests effectués. Les captures d'écran sont en \textbf{Annexe~D}.

% ──────────────────────────────────────────────────────────────────────────────
\section{Environnement de développement}

Le tableau~\ref{tab:env-dev} récapitule les outils constituant l'environnement de développement utilisé pour la réalisation du projet :

\begin{table}[H]
\centering
\begin{tabularx}{\textwidth}{llX}
\toprule
\rowcolor{eduai-primary}
\thead{Outil} & \thead{Version} & \thead{Usage} \\
\midrule
VS Code       & 1.90+    & Éditeur de code principal \\
Python        & 3.11     & Backend et services NLP \\
Node.js       & 20.17+   & Serveur frontend (Next.js) \\
Docker        & 24+      & Conteneurisation des services \\
Docker Compose& 2.x      & Orchestration multi-conteneurs \\
MongoDB Compass & 1.40+  & Administration de la base de données \\
Git           & 2.40+    & Gestion de versions \\
Postman       & 11+      & Tests manuels de l'API REST \\
\bottomrule
\end{tabularx}
\caption{Environnement de développement}
\label{tab:env-dev}
\end{table}

% ──────────────────────────────────────────────────────────────────────────────
\section{Technologies et frameworks}

Cette section présente les principales technologies et bibliothèques utilisées pour le développement de la plateforme.

\subsection{Backend — FastAPI (Python)}

Le tableau~\ref{tab:deps-backend} liste les principales dépendances du serveur backend :

\begin{table}[H]
\centering
\small
\begin{tabularx}{\textwidth}{llX}
\toprule
\rowcolor{eduai-primary}
\thead{Bibliothèque} & \thead{Version} & \thead{Rôle} \\
\midrule
FastAPI           & 0.115+ & Framework API REST asynchrone \\
Uvicorn           & 0.30+  & Serveur ASGI haute performance \\
Motor             & 3.6+   & Driver MongoDB asynchrone \\
Pydantic          & 2.x    & Validation et sérialisation des données \\
PyJWT             & 2.9+   & Génération et vérification JWT \\
Passlib[bcrypt]   & 1.7+   & Hachage sécurisé des mots de passe \\
PyMuPDF (fitz)    & 1.24+  & Extraction de texte des PDF \\
sentence-transformers & 3.0+ & Embeddings multilingues (MiniLM) \\
FAISS-cpu         & 1.8+   & Index vectoriel pour recherche sémantique \\
spaCy             & 3.7+   & NLP : tokenisation, segmentation \\
httpx             & 0.27+  & Client HTTP asynchrone (Groq API) \\
\bottomrule
\end{tabularx}
\caption{Principales dépendances backend}
\label{tab:deps-backend}
\end{table}

\subsection{Frontend — Next.js (React / TypeScript)}

Le tableau~\ref{tab:deps-frontend} présente les bibliothèques utilisées côté frontend :

\begin{table}[H]
\centering
\small
\begin{tabularx}{\textwidth}{llX}
\toprule
\rowcolor{eduai-primary}
\thead{Bibliothèque} & \thead{Version} & \thead{Rôle} \\
\midrule
Next.js          & 14.2    & Framework React avec SSR \\
React            & 18.3    & Bibliothèque UI \\
TypeScript       & 5.6     & Typage statique \\
Tailwind CSS     & 3.4     & Utilitaire CSS responsive \\
Zustand          & 4.5     & Gestion d'état léger \\
Axios            & 1.7     & Client HTTP \\
ReactFlow        & 11.11   & Visualisation de cartes mentales \\
Recharts         & 2.12    & Graphiques (activité, scores) \\
Framer Motion    & 11.11   & Animations et transitions \\
react-markdown   & 9.0     & Rendu Markdown des réponses LLM \\
react-dropzone   & 14.2    & Upload drag \& drop \\
\bottomrule
\end{tabularx}
\caption{Principales dépendances frontend}
\label{tab:deps-frontend}
\end{table}

\subsection{Services d'intelligence artificielle}

Le tableau~\ref{tab:services-ia} décrit les modèles et services d'IA intégrés dans le pipeline de traitement :

\begin{table}[H]
\centering
\begin{tabularx}{\textwidth}{lX}
\toprule
\rowcolor{eduai-primary}
\thead{Service} & \thead{Description} \\
\midrule
Groq Cloud (Llama 3.3 70B) & LLM principal : génération de résumés, quiz, flashcards, réponses Q\&A \\
paraphrase-multilingual-MiniLM-L12-v2 & Modèle d'embeddings (384 dimensions, support 50+ langues) \\
CamemBERT (fquad-piaf) & Modèle extractif de QA en français (fallback) \\
mT5 (XLSum) & Résumé extractif multilingue (fallback) \\
FAISS IndexFlatIP & Recherche par similarité cosinus (brute-force, exact) \\
\bottomrule
\end{tabularx}
\caption{Services d'intelligence artificielle}
\label{tab:services-ia}
\end{table}

% ──────────────────────────────────────────────────────────────────────────────
\section{Interfaces utilisateur}

La plateforme EduAI propose une interface web moderne et responsive, développée avec Next.js 14 et Tailwind CSS. L'interface comprend onze vues principales :

\begin{enumerate}
  \item \textbf{Page d'accueil} — présentation de la plateforme avec appel à l'action (CTA) ;
  \item \textbf{Authentification} — formulaires d'inscription et de connexion ;
  \item \textbf{Tableau de bord} — liste des cours de l'utilisateur sous forme de cartes ;
  \item \textbf{Upload} — interface de glisser-déposer pour l'import de fichiers PDF ;
  \item \textbf{Questions-Réponses} — chat conversationnel avec streaming des réponses ;
  \item \textbf{Résumés} — affichage structuré par chapitre avec concepts clés ;
  \item \textbf{Quiz} — QCM interactif avec calcul de score en temps réel ;
  \item \textbf{Flashcards} — cartes de révision avec évaluation de la qualité de réponse (SM-2) ;
  \item \textbf{Carte mentale} — visualisation interactive via ReactFlow ;
  \item \textbf{Examen} — mode examen chronométré avec configuration personnalisable ;
  \item \textbf{Administration} — panneau de gestion des utilisateurs et des statistiques.
\end{enumerate}

L'ensemble des captures d'écran illustrant ces interfaces est présenté en \textbf{Annexe~D}.

% ──────────────────────────────────────────────────────────────────────────────
\section{Tests et validation}

\subsection{Stratégie de test}

La validation du système repose sur plusieurs niveaux complémentaires, résumés dans le tableau~\ref{tab:tests} :

\begin{table}[H]
\centering
\begin{tabularx}{\textwidth}{llX}
\toprule
\rowcolor{eduai-primary}
\thead{Niveau} & \thead{Outil} & \thead{Couverture} \\
\midrule
Unitaire       & pytest        & Services (chunker, embedder, retriever) \\
Intégration    & pytest + httpx & Endpoints API complets \\
E2E LLM        & Scripts custom & Chaîne complète : upload → embed → Q\&A → résumé \\
Manuel         & Postman       & Scénarios utilisateur réels \\
Frontend       & Navigation    & Parcours utilisateur complets \\
\bottomrule
\end{tabularx}
\caption{Stratégie de test}
\label{tab:tests}
\end{table}

\subsection{Résultats des tests}

Le tableau~\ref{tab:test-results} présente les résultats de validation par module fonctionnel. L'ensemble des modules a été validé avec succès :

\begin{table}[H]
\centering
\begin{tabularx}{\textwidth}{lcX}
\toprule
\rowcolor{eduai-primary}
\thead{Module} & \thead{Statut} & \thead{Observations} \\
\midrule
Authentification     & \textcolor{green!60!black}{\checkmark} & JWT, bcrypt, rôles fonctionnels \\
Upload \& Indexation & \textcolor{green!60!black}{\checkmark} & PDF → chunks → FAISS opérationnel \\
Q\&A (RAG)           & \textcolor{green!60!black}{\checkmark} & Réponses pertinentes, streaming OK \\
Résumés              & \textcolor{green!60!black}{\checkmark} & Génération multi-chapitre, cache OK \\
Quiz                 & \textcolor{green!60!black}{\checkmark} & Génération, soumission, scoring OK \\
Flashcards SM-2      & \textcolor{green!60!black}{\checkmark} & Intervalles et ease\_factor corrects \\
Examen               & \textcolor{green!60!black}{\checkmark} & Timer, scoring, notes A-F OK \\
Admin                & \textcolor{green!60!black}{\checkmark} & CRUD utilisateurs, protection rôle \\
\bottomrule
\end{tabularx}
\caption{Résultats des tests par module}
\label{tab:test-results}
\end{table}
